\documentclass[11pt,letterpaper,oneside]{report}
\usepackage[T1]{fontenc}
\usepackage[utf8]{inputenc}
\usepackage[english]{babel}
\usepackage{lmodern}
\usepackage{microtype}
\usepackage[margin=0.86in,headheight=16pt]{geometry}
\usepackage[table,dvipsnames]{xcolor}
\usepackage{graphicx}
\usepackage{booktabs,tabularx,longtable,array}
\usepackage{enumitem}
\usepackage{titlesec}
\usepackage{fancyhdr}
\usepackage{hyperref}
\usepackage{xurl}
\Urlmuskip=0mu plus 1mu\relax
\usepackage{tcolorbox}
\usepackage{listings}
\usepackage{tikz}
\usetikzlibrary{arrows.meta,positioning,shapes.geometric,fit,calc}
\usepackage{needspace}
\usepackage{amssymb}
\definecolor{ink}{HTML}{14253A}
\definecolor{navy}{HTML}{17385D}
\definecolor{blue}{HTML}{20739A}
\definecolor{pale}{HTML}{ECF4F8}
\definecolor{green}{HTML}{1D6C5C}
\definecolor{softgreen}{HTML}{EEF8F4}
\definecolor{border}{HTML}{D4DFE7}
\definecolor{grey}{HTML}{506174}
\hypersetup{colorlinks=true,linkcolor=navy,urlcolor=blue,pdfauthor={Micro-Frontend Engineering Documentation Suite}}
\titleformat{\chapter}[display]{\normalfont\sffamily\bfseries\color{ink}}{\Large\color{blue}\MakeUppercase{\chaptername} \thechapter}{8pt}{\Huge}
\titlespacing*{\chapter}{0pt}{-10pt}{24pt}
\titleformat{\section}{\Large\sffamily\bfseries\color{navy}}{\thesection}{0.7em}{}
\titleformat{\subsection}{\large\sffamily\bfseries\color{navy}}{\thesubsection}{0.6em}{}
\setlist[itemize]{itemsep=3pt,topsep=5pt,leftmargin=1.6em}
\setlist[enumerate]{itemsep=4pt,topsep=5pt,leftmargin=1.8em}
\pagestyle{fancy}\fancyhf{}
\fancyhead[L]{\footnotesize\sffamily MICRO-FRONTEND ENGINEERING SUITE}
\fancyhead[R]{\footnotesize\sffamily\docshort}
\fancyfoot[L]{\footnotesize\color{grey}Engineering edition | October 2026}
\fancyfoot[R]{\footnotesize\thepage}
\renewcommand{\headrulewidth}{0.35pt}
\renewcommand{\arraystretch}{1.2}
\setlength{\emergencystretch}{2em}
\lstdefinestyle{eng}{basicstyle=\ttfamily\footnotesize,keywordstyle=\color{navy}\bfseries,commentstyle=\color{grey},stringstyle=\color{green},breaklines=true,breakatwhitespace=false,columns=fullflexible,keepspaces=true,showstringspaces=false,frame=single,rulecolor=\color{border},backgroundcolor=\color{pale},xleftmargin=5pt,xrightmargin=5pt,aboveskip=10pt,belowskip=10pt,captionpos=b}
\lstdefinelanguage{JavaScript}{morekeywords={const,let,var,require,return,async,await,import,export,default,function,try,catch,new,true,false,undefined,module},sensitive=true,morecomment=[l]//,morestring=[b]',morestring=[b]"}
\lstset{style=eng}
\newtcolorbox{sourcebox}{colback=pale,colframe=border,arc=2pt,boxrule=0.6pt,left=8pt,right=8pt,top=7pt,bottom=7pt}
\newtcolorbox{engineerbox}{colback=softgreen,colframe=green!45,arc=2pt,boxrule=0.6pt,left=8pt,right=8pt,top=7pt,bottom=7pt}
\newcommand{\based}[1]{\begin{sourcebox}\textbf{BOOK-GROUNDED CONCEPT}\\#1\end{sourcebox}}
\newcommand{\added}[1]{\begin{engineerbox}\textbf{ENGINEERING EXTENSION}\\#1\end{engineerbox}}
\newcommand{\term}[1]{\textbf{\color{navy}#1}}
\newcommand{\criterion}[1]{\textbf{Acceptance:} #1}
\newcommand{\book}{\emph{Building Micro-Frontends}, 2nd ed.}
\newcommand{\refs}{\chapter*{References and Provenance}\addcontentsline{toc}{chapter}{References and Provenance}
\begin{enumerate}[leftmargin=*]
\item \textbf{Primary source:} Luca Mezzalira, \emph{Building Micro-Frontends: Distributed Systems for the Frontend}, 2nd ed., O'Reilly Media, published November 2025 (copyright 2026), ISBN 978-1-098-17078-3. Chapter and printed-page references are included throughout. This suite paraphrases and synthesizes concepts; it does not reproduce the book.
\item \textbf{Engineering extension:} Webpack Module Federation reference: \url{https://webpack.js.org/plugins/module-federation-plugin}.
\item \textbf{Engineering extension:} Next.js multi-zones: \url{https://nextjs.org/docs/app/guides/multi-zones}. Details are version-sensitive and must be checked against the deployed framework version.
\item \textbf{Engineering extension:} OWASP Application Security Verification Standard 5.0: \url{https://owasp.org/www-project-application-security-verification-standard/}.
\item \textbf{Engineering extension:} GitHub Actions secure use reference: \url{https://docs.github.com/en/actions/reference/security/secure-use}.
\end{enumerate}}

\titleclass{\chapter}{straight}\titleformat{\chapter}[hang]{\huge\sffamily\bfseries\color{navy}}{\thechapter}{.7em}{}\titlespacing*{\chapter}{0pt}{24pt}{13pt}
\newcommand{\docshort}{STUDY GUIDES}
\begin{document}
\begin{titlepage}\thispagestyle{empty}
\vspace*{1.1in}{\Huge\sffamily\bfseries\color{navy} Chapter-by-Chapter Study Guides\par}
\vspace{18pt}{\Large\sffamily\color{blue} Fourteen guided readings, retrieval questions, applications, and assessments\par}
\vspace{27pt}\noindent\color{blue}\rule{0.75\textwidth}{1.4pt}
\vspace{23pt}\par{\large\sffamily A companion to Luca Mezzalira's \emph{Building Micro-Frontends}, Second Edition}\par
\vfill\noindent\textbf{Document series:} Micro-Frontend Software Architecture and Engineering\\
\textbf{Version:} 1.0 \quad \textbf{Prepared:} October 2026\\
\textbf{Production format:} pdfLaTeX-compatible standalone source\\
\textbf{Scope:} Original synthesis plus explicitly identified engineering extensions
\end{titlepage}
\pagenumbering{roman}\tableofcontents\clearpage\pagenumbering{arabic}
\chapter*{How to Use the Study Guides}\addcontentsline{toc}{chapter}{How to Use the Study Guides}
\based{The following 14 guides retain the second edition's chapter titles and ordering; page ranges refer to the book's printed pages. Each guide offers original learning objectives, review questions, and practical evidence, rather than reproducing the chapter text.}
\section*{Study protocol}
Before each chapter, define a hypothesis about business boundaries or system behavior. During reading, annotate the four architecture decisions affected. After reading, answer the questions without notes; then implement the exercise and retain the evidence. Assign each answer one of three states: \emph{understood}, \emph{uncertain}, or \emph{experiment required}. Revisit uncertain points after the lab.
\section*{Evaluation rubric}
\begin{tabularx}{\linewidth}{p{.18\linewidth}X}\toprule Level & Evidence standard \\\midrule
Foundational & Correctly defines terms and identifies the boundary or trade-off.\\
Practitioner & Produces an architecture/implementation with measurable integration tests.\\
Advanced & Defends the choice with alternatives, failure modes, rollback, and team ownership.\\\bottomrule\end{tabularx}

\chapter{Micro-Frontend Principles}
\based{Textbook alignment: Chapter 1, printed pp. 1--15.}
\section{Learning outcome}
Map autonomy, encapsulation, independent deployment, isolation, automation, and observability to frontend responsibilities.
\section{Reading checkpoints}
\begin{itemize}\item List the seven microservices-inspired principles and state a corresponding frontend test.
\item Explain why a young product can rationally begin as a frontend monolith.
\item Compare release autonomy with runtime isolation; they are not equivalent.
\item Identify the source of organizational coupling in a shared SPA.
\end{itemize}
\section{Retrieval and design questions}
Explain why a component library is not the same as autonomous delivery; identify two cases favoring a monolith. Give a counterexample in which the chapter's recommended direction would be inappropriate and explicitly identify the organizational constraint.
\paragraph{Architecture-review scenario} An executive wants micro-frontends for a three-person team: construct a no-decomposition argument and measurable revisit trigger.
\paragraph{Common analytical error} Assuming distributed frontends always outperform a well-modularized monolith.
\section{Applied assignment}
Choose a product and draw domain slices with ownership and release boundaries.
\section{Submission and grading}
\criterion{Capability map, one-page adoption hypothesis, baseline delivery metrics. Record the business reason, architectural constraint, primary risk, test method, and one credible alternative.}
\begin{engineerbox}\textbf{Engineering extension}\newline Treat the assignment as an architecture-review artifact: measure baseline behavior, preserve a reproducible test or configuration, and describe who owns post-deployment support.\end{engineerbox}

\chapter{Micro-Frontend Architectures and Challenges}
\based{Textbook alignment: Chapter 2, printed pp. 17--37.}
\section{Learning outcome}
Apply the four decisions: definition, composition, routing, communication. Introduce bounded contexts and distinguish horizontal from vertical cuts.
\section{Reading checkpoints}
\begin{itemize}\item Draw vertical and horizontal splits of the same page.
\item Identify a bounded context using workflows and ownership, not team preferences alone.
\item Select the composition layer and describe who resolves it.
\item Name routing, event, and API contracts separately.
\end{itemize}
\section{Retrieval and design questions}
What would invalidate your selected bounded context? How will domain events be governed? Give a counterexample in which the chapter's recommended direction would be inappropriate and explicitly identify the organizational constraint.
\paragraph{Architecture-review scenario} A checkout summary and cart appear together: justify whether they share a domain or communicate by contract.
\paragraph{Common analytical error} Treating visual placement as the sole input to domain boundaries.
\section{Applied assignment}
Record all four decisions as concise ADRs for a multi-team storefront.
\section{Submission and grading}
\criterion{Four ADRs, interface inventory, communication contracts. Record the business reason, architectural constraint, primary risk, test method, and one credible alternative.}
\begin{engineerbox}\textbf{Engineering extension}\newline Treat the assignment as an architecture-review artifact: measure baseline behavior, preserve a reproducible test or configuration, and describe who owns post-deployment support.\end{engineerbox}

\chapter{Discovering Micro-Frontend Architectures}
\based{Textbook alignment: Chapter 3, printed pp. 39--113.}
\section{Learning outcome}
Compare vertical and horizontal designs, client/server/edge composition, shell, Web Components, iframes, and architecture characteristics.
\section{Reading checkpoints}
\begin{itemize}\item Compare shell responsibilities against domain responsibilities.
\item Record browser, server, edge, Web Component, and iframe trade-offs.
\item Explain why combined horizontal views enlarge the compatibility matrix.
\item Capture a quality-attribute scenario and trade-off.
\end{itemize}
\section{Retrieval and design questions}
What extra runtime failure modes appear in a horizontal design? When does an iframe make sense? Give a counterexample in which the chapter's recommended direction would be inappropriate and explicitly identify the organizational constraint.
\paragraph{Architecture-review scenario} A route needs SEO but its personalization is sensitive: argue for a feasible composition model and cache boundary.
\paragraph{Common analytical error} Conflating independence of source repositories with independence of releases.
\section{Applied assignment}
Score three architectural candidates against explicit quality attributes.
\section{Submission and grading}
\criterion{Weighted trade-off matrix, architecture sketches, decision rationale. Record the business reason, architectural constraint, primary risk, test method, and one credible alternative.}
\begin{engineerbox}\textbf{Engineering extension}\newline Treat the assignment as an architecture-review artifact: measure baseline behavior, preserve a reproducible test or configuration, and describe who owns post-deployment support.\end{engineerbox}

\chapter{Client-Side Rendering Micro-Frontends}
\based{Textbook alignment: Chapter 4, printed pp. 115--149.}
\section{Learning outcome}
Work through a Module Federation host and domain remotes; study dependency sharing, caching, and independent releases.
\section{Reading checkpoints}
\begin{itemize}\item Map host and remote build artifacts and loading order.
\item Record dependencies to share versus duplicate.
\item Check version skew, CDN cache and remote unavailability.
\item Explain how legacy and new remotes can coexist during an incremental release.
\end{itemize}
\section{Retrieval and design questions}
What happens when a remote entry is unavailable or incompatible? Give a counterexample in which the chapter's recommended direction would be inappropriate and explicitly identify the organizational constraint.
\paragraph{Architecture-review scenario} The host and remote expect different singleton runtime versions: propose an experiment and safe fallback.
\paragraph{Common analytical error} Treating shared singleton dependencies as a free optimization.
\section{Applied assignment}
Create a shell loading a catalog remote plus a degraded-state fallback.
\section{Submission and grading}
\criterion{Running prototype, version contract, failure-injection evidence. Record the business reason, architectural constraint, primary risk, test method, and one credible alternative.}
\begin{engineerbox}\textbf{Engineering extension}\newline Treat the assignment as an architecture-review artifact: measure baseline behavior, preserve a reproducible test or configuration, and describe who owns post-deployment support.\end{engineerbox}

\chapter{Server-Side Rendering Micro-Frontends}
\based{Textbook alignment: Chapter 5, printed pp. 151--179.}
\section{Learning outcome}
Compare HTML fragments, route splitting, Next.js multi-zones, API strategies, SSR scale and caching layers.
\section{Reading checkpoints}
\begin{itemize}\item Differentiate HTML fragment assembly, route partitioning, and multi-zones.
\item Identify cache layers and personalization keys.
\item Analyze streaming, server scale, and initial response timing.
\item Trace session and authorization handling across zones.
\end{itemize}
\section{Retrieval and design questions}
Where should authentication execute? Which cache keys risk data leakage? Give a counterexample in which the chapter's recommended direction would be inappropriate and explicitly identify the organizational constraint.
\paragraph{Architecture-review scenario} Two zones use the same cookie; one leaks personalized HTML from the CDN: locate the boundary failure.
\paragraph{Common analytical error} Assuming SSR or CDN caching is safe for personalized content by default.
\section{Applied assignment}
Implement a reverse-proxy route split with a public and a private zone.
\section{Submission and grading}
\criterion{Route diagram, headers/cache plan, latency measurements. Record the business reason, architectural constraint, primary risk, test method, and one credible alternative.}
\begin{engineerbox}\textbf{Engineering extension}\newline Treat the assignment as an architecture-review artifact: measure baseline behavior, preserve a reproducible test or configuration, and describe who owns post-deployment support.\end{engineerbox}

\chapter{Micro-Frontend Automation}
\based{Textbook alignment: Chapter 6, printed pp. 181--205.}
\section{Learning outcome}
Study fast feedback, developer blueprints, version-control layouts, testing, fitness functions, and observability.
\section{Reading checkpoints}
\begin{itemize}\item Compare monorepo and polyrepo under independent artifact promotion.
\item Use a fitness function to guard a quality attribute automatically.
\item Separate component, contract, integration, and end-to-end tests.
\item Define developer blueprint and observability ownership.
\end{itemize}
\section{Retrieval and design questions}
How do you retain team autonomy with shared controls? What is a fitness function? Give a counterexample in which the chapter's recommended direction would be inappropriate and explicitly identify the organizational constraint.
\paragraph{Architecture-review scenario} A team can deploy independently but core E2E takes 80 minutes: redesign test layering without losing assurance.
\paragraph{Common analytical error} Using one full E2E suite as the only feedback loop.
\section{Applied assignment}
Introduce isolated contract and integration tests, plus a performance check.
\section{Submission and grading}
\criterion{CI pipeline, gate definitions, sample dashboards. Record the business reason, architectural constraint, primary risk, test method, and one credible alternative.}
\begin{engineerbox}\textbf{Engineering extension}\newline Treat the assignment as an architecture-review artifact: measure baseline behavior, preserve a reproducible test or configuration, and describe who owns post-deployment support.\end{engineerbox}

\chapter{Discover and Deploy Micro-Frontends}
\based{Textbook alignment: Chapter 7, printed pp. 207--224.}
\section{Learning outcome}
Examine frontend discovery, rollout configuration, canary/blue-green releases, and feature-flag boundaries.
\section{Reading checkpoints}
\begin{itemize}\item Model a frontend discovery manifest and its trust assumptions.
\item Distinguish blue-green, percentage canary and runtime feature flag.
\item Identify what a platform must change to promote without redeploying shell code.
\item Plan version/cohort consistency across a stateful journey.
\end{itemize}
\section{Retrieval and design questions}
Why separate release allocation from runtime feature toggles? Give a counterexample in which the chapter's recommended direction would be inappropriate and explicitly identify the organizational constraint.
\paragraph{Architecture-review scenario} A discovery record points to a release absent from CDN in one region: define rollout and rollback policy.
\paragraph{Common analytical error} Mistaking a runtime flag for release discovery and artifact integrity.
\section{Applied assignment}
Create a signed or trusted manifest registry for staged promotion.
\section{Submission and grading}
\criterion{Manifest schema, promotion history, rollback drill. Record the business reason, architectural constraint, primary risk, test method, and one credible alternative.}
\begin{engineerbox}\textbf{Engineering extension}\newline Treat the assignment as an architecture-review artifact: measure baseline behavior, preserve a reproducible test or configuration, and describe who owns post-deployment support.\end{engineerbox}

\chapter{Automation Pipeline for Micro-Frontends: A Case Study}
\based{Textbook alignment: Chapter 8, printed pp. 225--236.}
\section{Learning outcome}
Trace quality review, build, post-build inspection, and deployment across a domain-owned pipeline.
\section{Reading checkpoints}
\begin{itemize}\item Read the pipeline as distinct prebuild, build, postbuild and deployment stages.
\item Name the outputs and ownership of each quality gate.
\item Explain why deployment should promote immutable build output.
\item List failure response and recovery evidence.
\end{itemize}
\section{Retrieval and design questions}
What evidence should promotion verify rather than rebuild? Give a counterexample in which the chapter's recommended direction would be inappropriate and explicitly identify the organizational constraint.
\paragraph{Architecture-review scenario} CI scans pass but release manifest points to an unverified artifact: identify missing provenance control.
\paragraph{Common analytical error} Rebuilding independently on each deployment environment.
\section{Applied assignment}
Design a pipeline from commit to immutable artifact and staged release.
\section{Submission and grading}
\criterion{CI stages, artifact manifest, deployment audit record. Record the business reason, architectural constraint, primary risk, test method, and one credible alternative.}
\begin{engineerbox}\textbf{Engineering extension}\newline Treat the assignment as an architecture-review artifact: measure baseline behavior, preserve a reproducible test or configuration, and describe who owns post-deployment support.\end{engineerbox}

\chapter{Backend Patterns For Micro-Frontends}
\based{Textbook alignment: Chapter 9, printed pp. 237--268.}
\section{Learning outcome}
Compare service dictionary, gateway, domain API entry points, BFF, and GraphQL federation.
\section{Reading checkpoints}
\begin{itemize}\item Draw the request path for service dictionary, gateway, BFF and GraphQL.
\item Separate UI partitioning from authoritative domain data boundaries.
\item Write an API version and error compatibility policy.
\item Place authentication and resource-level authorization at trusted services.
\end{itemize}
\section{Retrieval and design questions}
Where is authorization authoritative? When is a separate BFF justified? Give a counterexample in which the chapter's recommended direction would be inappropriate and explicitly identify the organizational constraint.
\paragraph{Architecture-review scenario} Two domains request the same profile API: decide whether reuse belongs in BFFs, gateway, or shared client code.
\paragraph{Common analytical error} Allowing an API gateway or shell to become the shared business monolith.
\section{Applied assignment}
Implement a small BFF that aggregates two read-only API sources.
\section{Submission and grading}
\criterion{Sequence diagram, OpenAPI contract, backend security tests. Record the business reason, architectural constraint, primary risk, test method, and one credible alternative.}
\begin{engineerbox}\textbf{Engineering extension}\newline Treat the assignment as an architecture-review artifact: measure baseline behavior, preserve a reproducible test or configuration, and describe who owns post-deployment support.\end{engineerbox}

\chapter{Common Antipatterns in Micro-Frontend Implementations}
\based{Textbook alignment: Chapter 10, printed pp. 271--284.}
\section{Learning outcome}
Analyze misplaced boundaries, shared state, anarchy, anti-corruption layers, one-way sharing, and premature abstraction.
\section{Reading checkpoints}
\begin{itemize}\item Distinguish domain micro-frontend from common visual component.
\item Diagnose shared mutable state and bidirectional coupling.
\item Explain how anti-corruption layers aid migration.
\item Identify a premature abstraction and its lifecycle cost.
\end{itemize}
\section{Retrieval and design questions}
Which couplings force synchronized releases? How could an anti-corruption layer help? Give a counterexample in which the chapter's recommended direction would be inappropriate and explicitly identify the organizational constraint.
\paragraph{Architecture-review scenario} All micro-frontends import the same global store: propose incremental replacement without halting delivery.
\paragraph{Common analytical error} Creating many UI fragments that are not independently valuable business slices.
\section{Applied assignment}
Identify and refactor three intentional antipatterns in a sample architecture.
\section{Submission and grading}
\criterion{Smell register, before/after diagrams, dependency checks. Record the business reason, architectural constraint, primary risk, test method, and one credible alternative.}
\begin{engineerbox}\textbf{Engineering extension}\newline Treat the assignment as an architecture-review artifact: measure baseline behavior, preserve a reproducible test or configuration, and describe who owns post-deployment support.\end{engineerbox}

\chapter{Migrating to Micro-Frontends}
\based{Textbook alignment: Chapter 11, printed pp. 285--307.}
\section{Learning outcome}
Study iterative migration, prioritization, stakeholder engagement, UX, dependencies, auth and routing across old/new systems.
\section{Reading checkpoints}
\begin{itemize}\item Compare a big-bang rewrite with staged traffic cutover.
\item Select a first slice and state how you would measure success.
\item Plan shared auth, navigation and visual consistency across boundaries.
\item Define a stop/go criterion and rollback trigger.
\end{itemize}
\section{Retrieval and design questions}
Which first slice minimizes irreversibility? How do you avoid dual-login regressions? Give a counterexample in which the chapter's recommended direction would be inappropriate and explicitly identify the organizational constraint.
\paragraph{Architecture-review scenario} An extracted account domain still forces a joint release: trace coupling and refine the dependency plan.
\paragraph{Common analytical error} Extracting the hardest workflow first without proven rollback.
\section{Applied assignment}
Develop a strangler migration plan with cutover and rollback checkpoints.
\section{Submission and grading}
\criterion{Phased roadmap, risk register, sunset criteria. Record the business reason, architectural constraint, primary risk, test method, and one credible alternative.}
\begin{engineerbox}\textbf{Engineering extension}\newline Treat the assignment as an architecture-review artifact: measure baseline behavior, preserve a reproducible test or configuration, and describe who owns post-deployment support.\end{engineerbox}

\chapter{From Monolith to Micro-Frontends: A Case Study}
\based{Textbook alignment: Chapter 12, printed pp. 309--340.}
\section{Learning outcome}
Follow the fictional ACME Inc. decision process, vertical split, application shell, mounting, authentication, canaries, and design system.
\section{Reading checkpoints}
\begin{itemize}\item Identify the contextual reasons for ACME choosing a vertical split.
\item Trace shell initialization, lifecycle, routing and integration decisions.
\item Examine dependency and design-system choices during migration.
\item Review canary and localization issues without claiming a universal blueprint.
\end{itemize}
\section{Retrieval and design questions}
Which project constraints make a vertical split attractive? Which ACME choices are contextual? Give a counterexample in which the chapter's recommended direction would be inappropriate and explicitly identify the organizational constraint.
\paragraph{Architecture-review scenario} The ACME-inspired design performs poorly on mobile connections: assess whether to alter its chosen split.
\paragraph{Common analytical error} Copying case-study choices without validating local constraints.
\section{Applied assignment}
Reconstruct the decision record, then propose an alternative and defend it.
\section{Submission and grading}
\criterion{Case-study architecture review and comparison ADR. Record the business reason, architectural constraint, primary risk, test method, and one credible alternative.}
\begin{engineerbox}\textbf{Engineering extension}\newline Treat the assignment as an architecture-review artifact: measure baseline behavior, preserve a reproducible test or configuration, and describe who owns post-deployment support.\end{engineerbox}

\chapter{Introducing Micro-Frontends in Your Organization}
\based{Textbook alignment: Chapter 13, printed pp. 343--367.}
\section{Learning outcome}
Explore team topology, domain ownership, delivery data, RFC/ADR governance, and decentralized decisions.
\section{Reading checkpoints}
\begin{itemize}\item Contrast stream-aligned feature teams and central component teams.
\item Explain the intended purpose of an RFC versus an ADR.
\item Define a lightweight platform rule and a documented exception path.
\item Relate coordination burden to system boundaries and team topology.
\end{itemize}
\section{Retrieval and design questions}
When does a platform rule help rather than slow domain teams? Give a counterexample in which the chapter's recommended direction would be inappropriate and explicitly identify the organizational constraint.
\paragraph{Architecture-review scenario} Every team chooses its own auth library and telemetry format: establish minimal governance with autonomy.
\paragraph{Common analytical error} Governance by centralized approval of every domain-local implementation detail.
\section{Applied assignment}
Create an operating model with shared platform and domain responsibilities.
\section{Submission and grading}
\criterion{RACI, RFC template, scorecard, communication cadence. Record the business reason, architectural constraint, primary risk, test method, and one credible alternative.}
\begin{engineerbox}\textbf{Engineering extension}\newline Treat the assignment as an architecture-review artifact: measure baseline behavior, preserve a reproducible test or configuration, and describe who owns post-deployment support.\end{engineerbox}

\chapter{AI and Micro-Frontends: Augmenting, Not Replacing}
\based{Textbook alignment: Chapter 14, printed pp. 369--397.}
\section{Learning outcome}
Evaluate context engineering, tests, scaffolding, debugging, diagrams, migrations, code review, and human validation.
\section{Reading checkpoints}
\begin{itemize}\item Identify development tasks where AI can reduce toil without owning architecture judgment.
\item Design a context packet containing domain map, constraints and validation tests.
\item State how to independently validate generated tests, diagrams and migrations.
\item Explain hallucination, incomplete context and code-review risks.
\end{itemize}
\section{Retrieval and design questions}
What context is missing for a safe AI proposal? What must reviewers independently verify? Give a counterexample in which the chapter's recommended direction would be inappropriate and explicitly identify the organizational constraint.
\paragraph{Architecture-review scenario} An AI assistant suggests converting all routes to federation: define the human review and counterevidence needed.
\paragraph{Common analytical error} Accepting AI-generated architecture as verified simply because the code compiles.
\section{Applied assignment}
Use AI to draft a module contract, then critique it with tests and an ADR.
\section{Submission and grading}
\criterion{Prompt/context packet, review log, verified test results. Record the business reason, architectural constraint, primary risk, test method, and one credible alternative.}
\begin{engineerbox}\textbf{Engineering extension}\newline Treat the assignment as an architecture-review artifact: measure baseline behavior, preserve a reproducible test or configuration, and describe who owns post-deployment support.\end{engineerbox}

\chapter*{Final Synthesis Assessment}\addcontentsline{toc}{chapter}{Final Synthesis Assessment}
\section*{Capstone prompt}
An organization has eight frontend teams, a shared SPA, independently deployed backend services, and quarterly releases blocked by cross-domain regression testing. Design a business-domain decomposition and justify whether the first increment should use vertical routes, horizontal runtime composition, or no decomposition.
\section*{Scoring (100 points)}
\begin{itemize}\item Business-domain boundary and ownership (20).
\item Four-decision architecture record plus alternative (20).
\item Testable quality attributes and performance budget (15).
\item Security, auth, and API boundary design (15).
\item Continuous delivery, discovery, rollback, and telemetry (20).
\item Explicit trade-offs and migration sequence (10).\end{itemize}
\noindent Passing suggestion: at least 80/100, no missing security or rollback critical element. This is an instructional benchmark, not a requirement from the primary source.

\refs
\end{document}
